\chapter{Sinus, cosinus et oscillations}
\section*{Exercice 17.1 -- Angles et signes}
Par périodicité et symétrie, les valeurs sont $1/2$ et $-1/2$. L'équation a exactement les solutions $-\pi/3$ et $\pi/3$ sur l'intervalle demandé. Sur $\R$, il faudrait ajouter les multiples de $2\pi$.

\section*{Exercice 17.2 -- Inéquation}
Le cosinus est pair et décroît de $1$ à $-1$ sur $[0,\pi]$. Il atteint $-1/2$ en $2\pi/3$. L'ensemble des solutions est donc $[-2\pi/3,2\pi/3]$, bornes incluses.

\section*{Exercice 17.3 -- Composition}
La première dérivée est $6x\cos(3x^2)$. Pour la seconde fonction, $f'(x)=-2\sin(2x)$, négative sur $]0,\pi/2[$ puis positive sur $]\pi/2,\pi[$. Elle descend de $2$ à $0$ puis remonte à $2$.

\section*{Exercice 17.4 -- Rectangle inscrit}
L'aire est $4R^2\cos t\sin t=2R^2\sin(2t)$. Elle est au plus $2R^2$, atteinte pour $2t=\pi/2$, donc $t=\pi/4$. Les deux côtés sont alors égaux : le rectangle optimal est un carré de côté $R\sqrt2$. L'aire a bien l'unité d'une longueur au carré.

\section*{Exercice 17.5 -- Limite trigonométrique}
La première limite est $5\lim\sin(5x)/(5x)=5$. Pour la seconde, $1-\cos x=2\sin^2(x/2)$, donc le quotient vaut $\frac12[\sin(x/2)/(x/2)]^2\to1/2$. Les formules s'appliquent pour $x\ne0$ puis on passe à la limite.

\section*{Exercice 17.6 -- Tangente et domaine}
$\tan x=\sin x/\cos x$ est définie hors des $\pi/2+k\pi$. Sa dérivée vaut $1/\cos^2x>0$ sur chaque intervalle de définition. Le domaine n'est pas un intervalle : par exemple $\pi/4<3\pi/4$ mais les tangentes valent $1>-1$. La croissance est une propriété sur chaque composante.
